\SourceLeaf{356}%
rallelkreise, noch so geringfügig, dass man ihre Berücksichtigung bei den meisten Messungen kaum der wenn auch leichten Mühe werth halten wird.

In der Abhandlung ist die Theorie aller dieser und anderer damit zusammenhängenden Rechnungen vollständig entwickelt, an einer durchgehenden Musterrechnung erläutert, und mit einer Hülfstafel begleitet, die allerdings zunächst für diejenige Zone bestimmt ist, in welcher das Hannoversche Dreieckssystem liegt, aber auch ohne weiteres für Messungen benutzt werden kann, die diese Zone weit überschreiten: sie erstreckt sich nemlich über eine Zone von zwölf Breitengraden, in deren Mitte der gewählte Normal-Parallelkreis von \(52^0\,40'\) Breite liegt. Diese Tafel ist mit einer Schärfe berechnet, die ausreicht selbst wenn ein Dreieckssystem mit zehnziffrigen Logarithmen berechnet werden soll, also mit einer viel grösseren Schärfe, als man in den meisten Fällen beibehalten wird: indessen schien die kleine Raumersparniss, die durch Weglassung von ein paar Decimalen gewonnen sein würde, zu unerheblich, um beim Abdruck etwas davon zu unterdrücken.

Merklich und unmerklich sind bei Rechnungsoperationen relative Begriffe, und es ist also wohl der Mühe werth, sie nach ein paar aus der Abhandlung entlehnten Beispielen auf ein bestimmtes Maass zurückzuführen.

In dem Hannoverschen Dreieckssysteme ist das grösste Dreieck, welches auch zugleich am weitesten von dem Normal-Parallelkreise abliegt, dasjenige, welches zwischen den Punkten Brocken, Hohenhagen, Inselsberg gebildet wird. In diesem kommen daher auch die grössten Werthe der Richtungsreductionen vor, und zwar bei der Seite Hohenhagen-Inselsberg, wo die Reduction des Azimuths an dem erstern Endpunkte \(-0''00332\), am andern \(+0''00428\) beträgt. In dem ganzen Systeme kommen nur noch zwei andere Dreiecksseiten vor, wo die Reductionen \(0''001\) übersteigen, bei allen übrigen bleiben sie unter dieser Grösse.

Das grösste Hauptdreieck der trigonometrischen Vermessungen der Schweiz ist das zwischen den Punkten Chasseral, Suchet, Berra enthaltene; es berührt eben die südliche Grenze, bis zu welcher die Hülfstafel sich erstreckt, so dass die Richtungsreductionen sich noch vermittelst derselben berechnen lassen. Die grösste Reduction ist die, welche das Azimuth von Chasseral in Suchet trifft, und beträgt \(+0''06221\).

Es ist hieraus ersichtlich, dass in der ganzen Zone, worin das Hannoversche Dreieckssystem liegt, die Reduction ganz wegfällt, wenn die Rechnung auf
